\subsection{\texorpdfstring{第~\ref{sec:symmetry-and-reduction}~节}{对称性与约化}}\label{subsec:solutions-symmetry-and-reduction}

\begin{description}[font=\normalfont,leftmargin=2em,labelwidth=1.5em]
\RGsolution{ex:canonical-brackets-and-symplectic-coordinates}{canonical-brackets-and-symplectic-coordinates}
\begin{enumerate}[label=(\alph*),ref=\alph*,leftmargin=*]
\item 若 $\omega=\sum_{i=1}^n\mathrm{d}p_i\wedge\mathrm{d}x^i$, 则由 Exercise~\ref{ex:hamiltonian-mechanics}(\ref{ex:hamiltonian-fields-of-canonical-coordinates}) 知, 对 $i,j=1,\ldots,n$, 有 $\{x^i,x^j\}=\{p_i,p_j\}=0$ 和 $\{p_i,x^j\}=\delta_{ij}$. 另一方面, 若后一条件成立, 则
\[
  X_{x^i}=\sum_{j=1}^n\left(\{x^i,x^j\}\frac{\partial}{\partial x^j}
                  +\{x^i,p_j\}\frac{\partial}{\partial p_j}\right)
               =-\frac{\partial}{\partial p_i}
\]
及
\[
  X_{p_i}=\sum_{j=1}^n\left(\{p_i,x^j\}\frac{\partial}{\partial x^j}
                  +\{p_i,p_j\}\frac{\partial}{\partial p_j}\right)
               =\frac{\partial}{\partial x^i},
\]
这蕴含
\begin{align*}
  \omega\left(\frac{\partial}{\partial x^i},\frac{\partial}{\partial x^j}\right)
   &=\omega(X_{p_i},X_{p_j})=\{p_i,p_j\}=0;\\
  \omega\left(\frac{\partial}{\partial p_i},\frac{\partial}{\partial p_j}\right)
   &=\omega(X_{x^i},X_{x^j})=\{x^i,x^j\}=0;\\
  \omega\left(\frac{\partial}{\partial p_i},\frac{\partial}{\partial x^j}\right)
   &=-\omega(X_{x^i},X_{p_j})=-\{x^i,p_j\}=\delta_{ij},
\end{align*}
所以 $\omega=\sum_{i=1}^n\mathrm{d}p_i\wedge\mathrm{d}x^i$.
\item 由 Darboux 定理和 Exercise~\ref{ex:hamiltonian-mechanics}(\ref{ex:verify-canonical-symplectic-form}) 立即可知, $\omega^n=\omega\wedge\cdots\wedge\omega$ 是 $S$ 上的体积形式.
\item 设 $S$ 紧. 若 $\omega=\mathrm{d}\theta$, 则 $\omega^n=\mathrm{d}(\theta\wedge\omega\wedge\cdots\wedge\omega)$ (因为 $\mathrm{d}\omega=0$), 所以由 Stokes 定理,
\[
  \int_S\omega^n=\int_S\mathrm{d}(\theta\wedge\omega\wedge\cdots\wedge\omega)=0,
\]
这是矛盾.
\end{enumerate}

\RGsolution{ex:symplectic-form-from-nondegenerate-poisson-bivector}{symplectic-form-from-nondegenerate-poisson-bivector}
\begin{enumerate}[label=(\alph*),ref=\alph*,leftmargin=*]
\item 对 $i,j=1,\ldots,n$, 有
\[
  B^{ij}=B(\mathrm{d}x^i,\mathrm{d}x^j)=\{x^i,x^j\}.
\]
\item\label{sol:poisson-bivector-hamiltonian-field} 若 $F,G\in C^\infty(M)$, 则
\begin{align*}
  X_F\cdot G
  &=\{F,G\}=B(\mathrm{d}F,\mathrm{d}G)
   =\sum_{i,j=1}^n B^{ij}\frac{\partial F}{\partial x^i}\frac{\partial G}{\partial x^j}\\
  &=\left(\sum_{i,j=1}^n B^{ij}\frac{\partial F}{\partial x^i}
                                 \frac{\partial}{\partial x^j}\right)\cdot G.
\end{align*}
\item\label{sol:poisson-bivector-jacobi-identities} 对 $i,j,k=1,\ldots,n$, 由 Jacobi 恒等式, 有
\begin{align*}
  &\{x^i,\{x^j,x^k\}\}+\{x^j,\{x^k,x^i\}\}+\{x^k,\{x^i,x^j\}\}=0\\
  &\iff\{x^i,B^{jk}\}+\{x^j,B^{ki}\}+\{x^k,B^{ij}\}=0\\
  &\iff X_{x^i}\cdot B^{jk}+X_{x^j}\cdot B^{ki}+X_{x^k}\cdot B^{ij}=0.
\end{align*}
注意
\[
  X_{x^i}=\sum_{l=1}^n B^{il}\frac{\partial}{\partial x^l},
\]
得到
\[
  \sum_{l=1}^n\left(B^{il}\frac{\partial B^{jk}}{\partial x^l}
      +B^{jl}\frac{\partial B^{ki}}{\partial x^l}
      +B^{kl}\frac{\partial B^{ij}}{\partial x^l}\right)=0.
\]
\item 在辛流形 $(M,\omega)$ 上, Hamilton 向量场 $X_F$ 的定义可在局部坐标中写成
\[
  \iota(X_F)\omega=-\mathrm{d}F\iff
  \sum_{i,j=1}^n X_F^i\omega_{ij}\,\mathrm{d}x^i\otimes\mathrm{d}x^j
    =-\sum_{i=1}^n\frac{\partial F}{\partial x^i}\,\mathrm{d}x^i.
\]
利用 $(\omega_{ij})$ 的反对称性, 有
\[
  \sum_{j=1}^n\omega_{ij}X_F^j=\frac{\partial F}{\partial x^i}
     \iff X_F^i=\sum_{j=1}^n\omega^{ij}\frac{\partial F}{\partial x^j},
\]
其中 $(\omega^{ij}):=(\omega^{ij})^{-1}$. 因为由 (\ref{sol:poisson-bivector-hamiltonian-field}) 有
\[
  X_F^i=\sum_{j=1}^n B^{ji}\frac{\partial F}{\partial x^j}
         =-\sum_{j=1}^n B^{ij}\frac{\partial F}{\partial x^j},
\]
所以必有 $(B^{ij})=-(\omega^{ij})=-(\omega_{ij})^{-1}$.
\item 若 $B$ 非退化, 则对每个 $p\in M$, 可通过
\[
  \Phi(\omega)(\eta)=B(\omega,\eta)
\]
定义线性同构 $\Phi:T_p^*M\to T_pM$, 对所有 $\omega,\eta\in T_p^*M$ 成立. 在局部坐标中, 有
\[
  \Phi(\omega)(\eta)=\sum_{i,j=1}^n B^{ij}\omega_i\eta_j
      =\sum_{i,j=1}^n B^{ij}\omega_i\frac{\partial}{\partial x^j}(\eta),
\]
即
\[
  \Phi(\omega)=\sum_{i,j=1}^n B^{ij}\omega_i\frac{\partial}{\partial x^j}.
\]
于是可通过
\[
  \omega(v,w)=B(\Phi^{-1}(v),\Phi^{-1}(w))
\]
定义 2-形式 $\omega\in\Omega^2(M)$, 对所有 $v,w\in T_pM$ 成立. 令 $(B_{ij}):=(B^{ij})^{-1}$, 可在局部坐标下将 $\omega$ 写成
\begin{align*}
  \omega(v,w)
   &=\sum_{i,j,k,l=1}^n B^{ij}(-B_{ik}v^k)(-B_{jl}v^l)\\
   &=\sum_{k,l=1}^nB_{lk}v^kv^l=-\sum_{i,j=1}^n B_{ij}v^iv^j,
\end{align*}
也就是 $(\omega_{ij})=-(B_{ij})^{-1}$. 这说明 $\omega$ 确定与 $B$ 相同的 Poisson 括号, 且非退化. 为证明它是闭的, 注意到对于 $p,q,r=1,\ldots,n$, 有
\[
  \sum_{i,j,k,l=1}^n B_{pi}B_{qj}B_{rk}B^{il}\frac{\partial B^{jk}}{\partial x^l}
    =-\sum_{j,k=1}^n B_{qj}\frac{\partial B^{jk}}{\partial x^p}B_{kr}
    =\frac{\partial B_{qr}}{\partial x^p},
\]
所以用 $B_{pi}B_{qj}B_{rk}$ 乘 (\ref{sol:poisson-bivector-jacobi-identities}) 中的公式并对 $i,j,k$ 求和, 得
\[
  \frac{\partial B_{qr}}{\partial x^p}
      +\frac{\partial B_{rp}}{\partial x^q}
      +\frac{\partial B_{pq}}{\partial x^r}=0.
\]
这等价于
\[
  \sum_{i,j,k=1}^n\left(\frac{\partial\omega_{jk}}{\partial x^i}
       +\frac{\partial\omega_{ki}}{\partial x^j}
       +\frac{\partial\omega_{ij}}{\partial x^k}\right)
           \mathrm{d}x^i\otimes\mathrm{d}x^j\otimes\mathrm{d}x^k=0,
\]
又因为括号中的表达式对每一对指标反对称, 所以这又等价于
\[
  \sum_{i,j,k=1}^n\frac{\partial\omega_{jk}}{\partial x^i}
    \,\mathrm{d}x^i\wedge\mathrm{d}x^j\wedge\mathrm{d}x^k=0
    \iff\mathrm{d}\omega=0.
\]
\end{enumerate}

\RGsolution{ex:poisson-reduction-of-hyperbolic-geodesic-flow}{poisson-reduction-of-hyperbolic-geodesic-flow}
\begin{enumerate}[label=(\alph*),ref=\alph*,leftmargin=*]
\item 由群运算的表达式可知
\[
  (x,y)^{-1}=\left(-\frac xy,\frac1y\right),
\]
从而
\[
  L_{(a,b)^{-1}}(x,y)=\left(\frac xb-\frac ab,\frac yb\right).
\]
因此, 由 Example~\ref{ex:momentum-map-for-cotangent-lift}, $H$ 在自身上的作用到 $T^*H$ 的提升由
\[
  (a,b)\cdot(p_x\,\mathrm{d}x+p_y\,\mathrm{d}y)
   =(L_{(a,b)^{-1}})^*(p_x\,\mathrm{d}x+p_y\,\mathrm{d}y)
   =\frac{p_x}{b}\,\mathrm{d}x+\frac{p_y}{b}\,\mathrm{d}y
\]
给出, 在局部坐标下可写成
\[
  (a,b)\cdot(x,y,p_x,p_y)=\left(bx+a,by,\frac{p_x}{b},\frac{p_y}{b}\right).
\]
由于
\[
  K\left(bx+a,by,\frac{p_x}{b},\frac{p_y}{b}\right)
    =\frac{b^2y^2}{2}\left(\frac{p_x^2}{b^2}+\frac{p_y^2}{b^2}\right)
    =K(x,y,p_x,p_y),
\]
所以 $K$ 在 $H$ 下不变.
\item 函数 $F$ 和 $G$ 在 $H$ 下不变, 因为
\[
  F\left(bx+a,by,\frac{p_x}{b},\frac{p_y}{b}\right)
    =by\frac{p_x}{b}=yp_x=F(x,y,p_x,p_y)
\]
和
\[
  G\left(bx+a,by,\frac{p_x}{b},\frac{p_y}{b}\right)
    =by\frac{p_y}{b}=yp_y=G(x,y,p_x,p_y).
\]
这些函数是商流形 $T^*H/H$ 上的坐标 (它们是相对于左不变基的分量), 所以商的 Poisson 结构由
\begin{align*}
  \{F,G\}=X_F\cdot G
  &=\frac{\partial F}{\partial p_x}\frac{\partial G}{\partial x}
    +\frac{\partial F}{\partial p_y}\frac{\partial G}{\partial y}
    -\frac{\partial F}{\partial x}\frac{\partial G}{\partial p_x}
    -\frac{\partial F}{\partial y}\frac{\partial G}{\partial p_y}\\
  &=-p_xy=-F
\end{align*}
确定. 因此商上的 Poisson 双向量为
\begin{align*}
  B&=\{F,G\}\frac{\partial}{\partial F}\otimes\frac{\partial}{\partial G}
      +\{G,F\}\frac{\partial}{\partial G}\otimes\frac{\partial}{\partial F}\\
   &=-F\frac{\partial}{\partial F}\otimes\frac{\partial}{\partial G}
      +F\frac{\partial}{\partial G}\otimes\frac{\partial}{\partial F}.
\end{align*}
因为 $F=0$ 时 $B$ 为零, 商 $T^*H/H$ 不是辛流形.
\item 沿经过单位元 $e=(0,1)$ 的曲线 $(a(t),b(t))$ 对表达式
\[
  L_{(a,b)}(x,y)=(bx+a,by)
\]
求导, 容易看出 $V=\alpha\frac{\partial}{\partial x}+\beta\frac{\partial}{\partial y}\in\mathfrak h$ 的无穷小作用为
\[
  X^V=(\alpha+\beta x)\frac{\partial}{\partial x}+\beta y\frac{\partial}{\partial y}.
\]
由 Example~\ref{ex:momentum-map-for-cotangent-lift}, $H$ 在 $T^*H$ 上作用的动量映射是由
\[
 J(p_x\,\mathrm{d}x+p_y\,\mathrm{d}y)(V)
   =(p_x\,\mathrm{d}x+p_y\,\mathrm{d}y)(X^V)
   =(\alpha+\beta x)p_x+\beta yp_y
\]
给出的映射 $J:T^*H\to\mathfrak h^*$. 因为 $K$ 在 $H$ 下不变, $J$ 沿 $K$ 的 Hamilton 流为常数, 所以取 $\alpha=0$ 和 $\beta=1$, 得到 $K$ 的 Hamilton 流的非平凡第一积分
\[
  I(x,y,p_x,p_y)=xp_x+yp_y
\]
(除了显然的第一积分 $K$ 和 $p_x$ 之外). 因此, 满足 $K=E$, $p_x=l$ 和 $I=m$ 的测地线满足
\[
  y^2(p_x^2+p_y^2)=2E\iff y^2l^2+(m-xl)^2=2E,
\]
当 $l\ne0$ 时, 这是圆心在 $x$ 轴上的圆的方程.
\end{enumerate}
\end{description}
